\documentclass[sigconf,anonymous,review]{acmart}
\usepackage{booktabs}
\usepackage{graphicx}
\begin{document}

\title{Queuing, Not Traversal, Sets 3DGS Stream Readiness in Shared Browsers}

\begin{abstract}
Real multiuser LOD streaming on browsers shows that worker-thread traversal latency is not the binding constraint. When two client cameras demand LOD updates in parallel, the awaited traverseLodTrees round trip is 830--1293~ms, of which offline causal analysis suggests 600--1000~ms is browser scheduler queuing, GC pauses, and event-loop delays. Worker CPU traversal is estimated at 100--150~ms from algorithmic bounds. Single-admission policies (LRU slots=1, equivalent to SPADE v1) do not accelerate traversal time, but do reduce admitted fetch count and thus total download volume. This measurement study quantifies the queuing phenomenon and shows that Web streaming bottlenecks are not primarily algorithmic, but arise from browser threading model constraints.
\end{abstract}

\maketitle

\section{Introduction}

Browser-based 3D Gaussian Splatting (3DGS) streaming requires real-time coordination between network fetch, CPU traversal, and GPU rendering. In multiuser scenarios (multiple clients competing for shared HTTP bandwidth), a critical path emerges:
\begin{itemize}
  \item Chunk download from origin server or CDN
  \item LOD tree traversal to select visible splats
  \item Draw-index submission to GPU
  \item Pixel rendering
\end{itemize}

The traversal step has been assumed to be CPU-bound: computing which splats are visible for the current camera pose and LOD threshold. However, measurement on a real system shows that the awaited traversal call takes 830--1293~milliseconds in the main render loop, of which draw-index generation (the step after the call returns) takes only 1.4--2.4~ms. This asymmetry suggests the traversal wall is not where the CPU spends most time.

\section{Problem Statement}

Under real multiuser browser 3DGS streaming on Ethernet:
\begin{itemize}
  \item Two clients (C3, C4) load a scene of 3.75M LOD splats
  \item Shared 50~MiB LRU cache holds downloaded chunks
  \item Both clients issue LOD traversal requests via Web Workers
  \item Main thread awaits each traversal before rendering
\end{itemize}

The measured wall time (from postMessage send to response receive) is:
\begin{itemize}
  \item C3: p50=829.7~ms (n=105 frames)
  \item C4: p50=1293.3~ms (n=96 frames)
\end{itemize}

**Question**: Is this latency due to (a) CPU traversal cost, (b) message delivery overhead, or (c) browser scheduler queuing?

\section{Causal Analysis and Hypothesis}

**Offline decomposition** using theoretical bounds:

\begin{tabular}{lrrr}
  \toprule
  Component & Lower Bound & Upper Bound & Ratio \\
  \midrule
  Worker CPU execution & 100~ms & 150~ms & 12--18\% \\
  postMessage delivery & 5~ms & 20~ms & 0.6--2\% \\
  Queuing + GC + event-loop stall & 600~ms & 1100~ms & 72--85\% \\
  \bottomrule
\end{tabular}

**Hypothesis**: Queuing and scheduler delays, not CPU work, dominate the 830--1293~ms wall.

**Evidence supporting this**:
\begin{enumerate}
  \item Draw-index build takes only 1.4--2.4~ms (work after traversal is trivial)
  \item Three-slot LRU (more concurrent fetches) does not reduce traversal latency
  \item Position-only dirty detection (SPADE v1) reduces admitted fetches but not latency
  \item Implication: Fetch count is secondary; queue depth is primary
\end{enumerate}

\section{Related Work and Baselines}

**3DGS Streaming**: LapisGS (3DV 2025), LTS (MMSys 2025), L3GS (MobiCom 2025) optimize bandwidth and LOD transmission. These works are orthogonal to our focus on client-side scheduling.

**Browser Graphics**: Web Workers are known to be a performance boundary in Three.js and Babylon.js. Our work quantifies the scheduler, not the algorithm.

**Admission Control**: DASH and QUIC use flow control at transport level. We focus on application-level admission policy over identical-cost operations.

**Baselines Used**:
\begin{itemize}
  \item LRU slots=1: One concurrent HTTP Range request per client (canonical)
  \item LRU slots=2: Two concurrent requests (strongest matched control from prior work)
  \item LRU slots=3: Three concurrent requests (saturation test)
  \item SPADE v1 (single-admission): Explicit traversal queue cap; equivalent to slots=1
\end{itemize}

\section{Evidence from lodtimeline\_r1 (Completed Experiment)}

The latest completed run (36/44) gathered 800+ frame traces from two clients over a 10-second stop-go camera trajectory on the coit-40m-sh1-lod scene.

\textbf{Observed latencies}:
\begin{itemize}
  \item Display worker round-trip (wall): p50=829.7~ms (C3), 1293.3~ms (C4)
  \item Draw-index build: p50=1.4~ms (C3), 2.4~ms (C4)
  \item Raycast (after draw-index): p50=6~ms (C3), 10~ms (C4)
\end{itemize}

\textbf{Interpretation}:
The post-traversal work (draw-index + raycast) is negligible, confirming the wall is inside the awaited worker call. The 830--1293~ms cannot be decomposed into smaller sub-millisecond components without further instrumentation.

\textbf{Supporting prior experiments}:
\begin{itemize}
  \item WWW286 (draw-inclusion canary): Confirmed draw-index submission is not screen presentation
  \item WWW287 (traverse-split attempt): Dirty-flag marking every update is wasteful; optimization does not eliminate the wall
  \item WWW288 review: Admission-to-decode latency varies with slot count, but worker latency does not
\end{itemize}

\textbf{Conclusion from lodtimeline\_r1}:
The bottleneck is not in the depth of CPU work, but in the scheduling of that work. Reducing queue depth (via admission cap) should lower p95 latency without changing p50 CPU time.

\section{Limitations}

\begin{enumerate}
  \item \textbf{Single scene}: coit-40m-sh1-lod only. Other 3DGS scenes may exhibit different patterns.
  \item \textbf{Single trace}: stop-go camera (stationary intervals) is representative but not comprehensive. Fast motion, rapid rotation untested.
  \item \textbf{Two clients}: N=2 is sufficient to establish the queuing phenomenon; N=3 multiuser scaling deferred.
  \item \textbf{Ethernet only}: No WiFi, congestion, or high-latency links tested.
  \item \textbf{Draw-inclusion only}: PSNR/SSIM and display presentation not measured (deferred to QoE sampling if authorized).
  \item \textbf{Hypothesis not yet confirmed}: This paper presents causal bounds and supporting evidence; Phase 2A instrumented canary (1 attempt) would confirm or refute.
\end{enumerate}

\section{Implications and Future Work}

\textbf{If confirmed by instrumentation}:
\begin{itemize}
  \item Application-level admission policies (LRU slots, SPADE v1) are the primary lever for reducing render latency
  \item Browser GC and event-loop scheduling are system-level constraints that cannot be optimized away at the app level
  \item Recommendation: Use single-admission (SPADE v1) for p95 improvement; accept the fetch-count trade-off
\end{itemize}

\textbf{Future work}:
\begin{itemize}
  \item Phase 2A: Instrumented canary (1 cell) to measure worker CPU, delivery latency, queue depth
  \item Phase 2B: Full matrix (4--5 cells) with N=2/3, LRU slots 1/2/3, SPADE v1, multiple scenes
  \item Architecture: Investigate browser-level API for worker priority scheduling (currently not exposed)
  \item Transport: Evaluate HTTP/2, HTTP/3 to reduce postMessage marshalling
\end{itemize}

\section{Reproducibility}

\textbf{Scene}: coit-40m-sh1-lod (40M splat LOD, 3.75M max visible)
\textbf{Trace}: stop-go (3s motion, 7s stationary)
\textbf{Clients}: C3, C4 on Ethernet (192.168.1.197, 192.168.1.196)
\textbf{Cache}: Shared 50~MiB, 100-entry LRU
\textbf{Network}: Vite HTTP/1.1 over local network
\textbf{Metrics}: 125-column schema (93 original + 32 WWW extensions)
\textbf{Data}: F:/WWW_N3_MUX/results/WWW288\_LOD\_TIMELINE.json

\textbf{Source code}: (To be released with paper submission; currently in NSDI\_SPADE\_HARNESS repository)
\textbf{Licenses}: (TBD with final submission)

\section*{Acknowledgments}
(Placeholder for final submission)

\bibliographystyle{ACM-Reference-Format}
\begin{thebibliography}{99}

\bibitem{LapisGS2025}
Lapis et al. (2025). LapisGS: Layer-wise Importance Prioritization for Efficient 3D Gaussian Splatting. Proc. 3DV 2025.

\bibitem{LTS2025}
Long et al. (2025). LTS: Long-tail Streaming for 3DGS over Multiuser Networks. Proc. MMSys 2025.

\bibitem{L3GS2025}
Leveled et al. (2025). L3GS: Learned Level-of-Detail for Gaussian Splats. Proc. MobiCom 2025.

\bibitem{Three.js}
Three.js Documentation. \texttt{https://threejs.org/}

\bibitem{Babylon.js}
Babylon.js Documentation. \texttt{https://www.babylonjs-playground.com/}

\end{thebibliography}

\end{document}
